\scanpage{251}
\mspage{250}
\mspage{2}\hfill n$^\circ$ 293

un corps des fractions de $o/\unsure{p|\mathfrak p}$ ; soit $R(\theta_1,\ldots,\theta_n)$ une extension transcendante pure de $R$ engendrée par $n$ éléments $\theta_1,\ldots,\theta_n$ algébriquement indépendants par rapport à $R$. L'homomorphisme canonique de $o$ sur $o/p$ se prolonge en un homomorphisme $\omega$ de $o[t_1,\ldots,t_n]$ dans $R[\theta_1,\ldots,\theta_n]$ qui applique $t_i$ sur $\theta_i$ $(1\leq i\leq n)$ ; le noyau de cet homomorphisme est $po[t_1,\ldots,t_n]$, comme il résulte tout de suite du fait que $\theta_1,\ldots,\theta_n$ sont algébriquement indépendants sur $R$ ; par ailleurs, $\omega(o[t_1,\ldots,t_n])$ est un domaine d'intégrité ce qui montre bien que $po[t_1,\ldots,t_n]$ est premier. De plus $\omega(o[t_1,\ldots,T_n])$ n'est pas un corps. En effet, dans le cas contraire, ce corps devrait contenir $R$, donc être identique à $R[\theta_1,\ldots,\theta_n]$ or ce dernier anneau n'est pas un corps puisque $n>0$. On en conclut que $po[t_1,\ldots,t_n]$ n'est pas un idéal premier maximal de $o[t_1,\ldots,t_n]$. Par contre, $q$ étant un idéal premier maximal de $r$, et $r$ étant entier sur $o[t_1,\ldots,t_n]$, $q\cap o[t_1,\ldots,t_n]$ est un idéal premier maximal de $o[t_1,\ldots,t_n]$ ; cet idéal contient $po[t_1,\ldots,t_n]$, et en est distinct en vertu de ce qu'on vient de dire. Comme $o[t_1,\ldots,t_n]$ est intégralement clos, on sait qu'il y a un idéal premier $\unsure{n|\mathfrak n}$ de $r$ tel que $n\subset q$, $n\cap o[t_1,\ldots,t_n]=po[t_1,\ldots,t_n]$ ; cet idéal est contenu dans $q$, distinct de $q$ et contient $p$.

\uline{Lemme} 2 - \uline{Soient} $o$ \uline{un domaine d'intégrité local,} $m$ \uline{son idéal premier maximal,} $t$ \uline{un élément d'un sur-corps du corps des fractions de} $o$ ; \uline{on suppose que} $o$ \uline{est intégralement fermé dans} $o[t]$, \uline{que} $t\notin o$, \uline{et que l'idéal engendré par} $m$ \uline{dans} $o[t]$ \uline{ne contient pas} 1. \uline{Si} $F$ \uline{est un polynôme à coefficients dans} $o$ \uline{tel}
